\documentclass[10pt,a4paper]{amsart}
\usepackage[margin=19mm]{geometry}
\usepackage{amsmath,amssymb,mathtools}
\usepackage[scr=rsfs,cal=euler]{mathalfa}
\usepackage{xfrac}
\usepackage[unicode,hidelinks,bookmarks=false]{hyperref}
\newcommand{\ie}{i.e.\ }
\newcommand{\cf}{cf.\ }
\newcommand{\resp}{respectively\ }
\newcommand{\Subsection}{\subsection}
\providecommand{\nobreakdash}{\nobreak\hbox{-}}
\setlength{\emergencystretch}{2em}
\multlinegap0pt
\usepackage{amssymb}
\usepackage{xcolor}
\newtheorem{proposition}{Proposition}
\newtheorem{conjecture}{Conjecture}
\newtheorem{theorem}{Theorem}
\title{Computational Evidence Against Quadratic--Cubic Factorization for the Second Cuboid Quintic}
\author{Valery Asiryan, Randall L. Rathbun}
\date{Source: arXiv:2601.07899v2 (CC BY 4.0)}
\begin{document}
\maketitle

\noindent\emph{Source and licence.} Computational Evidence Against Quadratic--Cubic Factorization for the Second Cuboid Quintic. Version arXiv:2601.07899v2 (CC BY 4.0), distributed by arXiv under the Creative Commons Attribution 4.0 International licence, http://creativecommons.org/licenses/by/4.0/, as recorded in the arXiv licence metadata for this exact version and shown on the abstract page https://arxiv.org/abs/2601.07899v2. Full licence notice: \texttt{licenses/arxiv-2601.07899-CC-BY-4.0.txt}.\par\smallskip

\noindent\textit{Editorial scope.} Counted unit: the article's theorem thm:main (source0.tex lines 395--399). The article's complete original proof of it is delivered with it (source0.tex lines 401--412, one \texttt{proof} environment of 12 lines). No mathematics is written, repaired or re-derived: the statement and the proof are the article's own lines, in the article's own notation, and no retained line is substituted or re-flowed.  Retained uncounted as the source-local context the proof needs: lines 278--290; lines 326--336; lines 377--379; lines 391--393.  Nothing was omitted from any retained span, so the package carries no omission record.  No retained line contains a citation, so the wrapper carries no bibliography and no bibliography entry.  No retained line contains a figure environment, an image-inclusion command or drawing code, so no image is delivered and the package declares no asset dependency.  Editorial formatting only, in addition: an AMS \texttt{amsart} preamble that restates the article's own declarations and macros used by the retained lines, namely the environments \texttt{proposition}, \texttt{conjecture}, \texttt{theorem} on separate counters, together with the standard packages those lines need.  The retained environments print in this document as Proposition 1, Conjecture 1, Theorem 1 in order of appearance, whereas the article prints the same items with the numbering of its own full document.  The complete native source of the article as distributed in the arXiv e-print for this exact version is bundled unchanged as \texttt{sources/192551/source0.tex}.
\begin{proposition}[Divisibility criterion]\label{prop:criterion}
Fix $s\in\mathbb{Q}$.
The following are equivalent:
\begin{enumerate}
\item $P_s(x)$ admits a $2+3$ factorization over $\mathbb{Q}$;
\item there exist $a,b\in\mathbb{Q}$ such that $R_1(s,a,b)=R_0(s,a,b)=0$.
\end{enumerate}
Moreover, if such $a,b$ exist then either
\begin{enumerate}
\item[(A)] $L(s,a)\neq 0$ and $F(s,a)=0$, in which case necessarily $b=C(s,a)/L(s,a)$; or
\item[(B)] $L(s,a)=0$ and $C(s,a)=0$ (the \emph{degenerate locus}).
\end{enumerate}
\end{proposition}

\begin{proposition}[Degenerate locus classification]\label{prop:degenerate}
The system
\[
L(s,a)=0,\qquad C(s,a)=0
\]
has the following rational solutions:
\[
(s,a)=(1,2)\quad\text{and}\quad (s,a)=(-1,2).
\]
In particular, for $s\in\mathbb{Q}_{>0}$ with $s\neq 1$ the degenerate locus is empty.
\end{proposition}

\begin{equation}\label{eq:CQ}
\{(1:2:1),\ (0:1:0),\ (0:0:1),\ (0:6:1),\ (1:0:0),\ (-1:2:1),\ (1:0:1),\ (-1:1:0)\}.
\end{equation}

\begin{conjecture}[Completeness of the rational-point list]\label{conj:complete}
The points listed in \eqref{eq:CQ} are all rational points on $\overline{C}$, i.e.\ $\overline{C}(\mathbb{Q})$ consists of exactly these $8$ points.
\end{conjecture}

\begin{theorem}[Conditional exclusion of $2+3$ factorization]\label{thm:main}
Assume Conjecture~\ref{conj:complete}.
Let $s\in\mathbb{Q}_{>0}$ with $s\neq 1$.
Then the quintic $P_s(x)$ admits no $2+3$ factorization over $\mathbb{Q}$ (equivalently, it has no quadratic factor over $\mathbb{Q}$).
\end{theorem}

\begin{proof}
Assume, for contradiction, that $P_s$ admits a $2+3$ factorization over $\mathbb{Q}$ for some $s\in\mathbb{Q}_{>0}$ with $s\neq 1$.
Then by Proposition~\ref{prop:criterion} there exist $a,b\in\mathbb{Q}$ with $R_1(s,a,b)=R_0(s,a,b)=0$.

By Proposition~\ref{prop:degenerate}, the degenerate case $L=C=0$ can occur for $s>0$ only when $s=1$, which is excluded.
Hence $L(s,a)\neq 0$ and Proposition~\ref{prop:criterion}(A) applies, giving $F(s,a)=0$.
Thus $(s,a)$ is an affine rational point on $\overline{C}$.

By Conjecture~\ref{conj:complete}, every affine rational point on $\overline{C}$ has $s\in\{-1,0,1\}$.
Since $s>0$ and $s\neq 1$, no such point exists.
This contradiction shows that $P_s$ has no quadratic factor over $\mathbb{Q}$, hence admits no $2+3$ factorization.
\end{proof}

\end{document}
